\SGASection{2}{Structure of the unipotent radical of a parabolic subgroup}
\SGARefTarget{sga3:xxvi:section:2}

\begin{sgastatement}{\SGAInlineRefTarget{sga3:xxvi:proposition:2-1}Proposition 2.1}
Let \(S\) be a scheme, \(G\) an \(S\)-reductive group, \(P\) a parabolic
subgroup of \(G\), and \(\rad^u(P)\) its unipotent radical.  There is a
sequence of sub-group-schemes of \(\rad^u(P)\)
\[
 \rad^u(P)=U_0\supset U_1\supset U_2\supset\cdots
 \supset U_n\supset\cdots
\]
having the following properties:
\begin{enumerate}[label=\textup{(\roman*)}]
\item Every \(U_i\) is smooth, with connected fibers, characteristic and
closed in \(P\).  The commutator of a section of \(U_i\) and a section of
\(U_j\) is a section of \(U_{i+j+1}\), over a varying \(S'\to S\).
\item For every \(i\geq0\), there is a locally free
\(\mathcal O_S\)-module \(\mathcal E_i\) and an isomorphism of sheaves of
\(S\)-groups
\[
 U_i/U_{i+1}\xrightarrow{\sim}W(\mathcal E_i).
\]
Moreover, the automorphisms of \(P\), over a varying \(S'\to S\), act
linearly on \(W(\mathcal E_i)\).
\item For every \(s\in S\), one has \(U_{n,s}=e\) whenever
\[
 n>\dim\bigl(\rad^u(P)_s\bigr).
\]
\end{enumerate}
\end{sgastatement}
\SGARefTarget{sga3:xxvi:prop:2:1}

\paragraph{\SGARefLink{sga3:xxvi:prop:2:1}{2.1.1}.}
Suppose first that the pair \((G,P)\) admits a pinning.  Let
\((T,M,R,\Delta,\ldots)\) be a pinning of \(G\) adapted to \(P\), and let
\(\Delta(P)\subset\Delta\) be determined by \(P\).  Write
\(\alpha_1,\ldots,\alpha_p\) for the elements of \(\Delta(P)\), and
\(\beta_1,\ldots,\beta_q\) for the elements of
\(\Delta-\Delta(P)\).  Every root \(\gamma\in R\) has a unique expression
\[
 \gamma
 =a_1\alpha_1+\cdots+a_p\alpha_p
  +b_1\beta_1+\cdots+b_q\beta_q.
\]
Set
\[
 a(\gamma)=b_1+\cdots+b_q.
\tag{7}\label{sga3:xxvi:proposition:2-1:equation:7}
\]
The definitions immediately give (cf.\ \SGARefLink{sga3:xxvi:proposition:1-12}{1.12}):
\[
\begin{aligned}
 U_\gamma\subset P
   &\ \Longleftrightarrow\ a(\gamma)\geq0,\\
 U_\gamma\subset\rad^u(P)
   &\ \Longleftrightarrow\ a(\gamma)>0,\\
 a(n\gamma+m\gamma')
   &=n\,a(\gamma)+m\,a(\gamma')
   \qquad(n,m\in\mathbb Z).
\end{aligned}
\]

For \(i\geq0\), let \(R_i\) be the set of roots \(\gamma\in R\) such that
\(a(\gamma)>i\).  Each \(R_i\) is a closed set of roots satisfying
\(R_i\cap(-R_i)=\varnothing\).  Consider (Exp.\ XXII, \SGARefLink{sga3:xxii:cor:5.6.5}{5.6.5}) the
\(S\)-group
\[
 U_i=U_{R_i}=\prod_{\gamma\in R_i}U_\gamma.
\]
It is a closed sub-group-scheme of \(G\), smooth over \(S\), with connected
fibers.

Let \(\alpha,\beta\in R\).  Consider the commutation relation of
Exp.\ XXII, \SGARefLink{sga3:xxii:cor:5.5.2}{5.5.2}:
\[
 p_\alpha(x)p_\beta(y)p_\alpha(-x)
 =
 p_\beta(y)
 \prod_{n,m\in\mathbb N^*}
 p_{n\alpha+m\beta}
 \bigl(C_{n,m,\alpha,\beta}x^ny^m\bigr),
\]
where every
\[
 p_\gamma:\mathbb G_{a,S}\xrightarrow{\sim}U_\gamma
\]
is an isomorphism of vector groups.  If \(a(\alpha)>i\) and
\(a(\beta)>j\), then, for \(n,m>0\),
\[
\begin{aligned}
 a(n\alpha+m\beta)
 &=n\,a(\alpha)+m\,a(\beta)\\
 &\geq n(i+1)+m(j+1)>i+j+1.
\end{aligned}
\]
Hence the commutator of a section of \(U_\alpha\) and a section of
\(U_\beta\) belongs to \(U_{i+j+1}\); therefore
\[
 (U_i,U_j)\subset U_{i+j+1}.
\]
In particular \(U_i/U_{i+1}\) is commutative, and identifies naturally
with
\[
 U_i/U_{i+1}
 \simeq
 \prod_{a(\gamma)=i+1}U_\gamma
 \simeq W(\mathcal E_i),
\]
where \(\mathcal E_i\) is the direct sum of the
\(\mathfrak g^\gamma\) for which \(a(\gamma)=i+1\).

Return to the commutation formula and suppose
\(a(\alpha)\geq0\), \(a(\beta)>i\).  If \(n,m\in\mathbb N^*\), then
either
\[
 a(n\alpha+m\beta)>i+1,
\]
or \(a(n\alpha+m\beta)=i+1\), in which case necessarily \(m=1\).
This first proves that \(\Int(p_\alpha(x))\) preserves \(U_i\), and hence
also \(U_{i+1}\).  It then shows that, modulo \(U_{i+1}\), the expression
for \(\Int(p_\alpha(x))p_\beta(y)\) contains only terms
\[
 p_{n\alpha+\beta}
 \bigl(C_{n,1,\alpha,\beta}x^ny\bigr),
\]
which are linear in \(y\).  Thus the inner automorphisms defined by
sections of \(U_\alpha\) act linearly on
\(U_i/U_{i+1}\simeq W(\mathcal E_i)\).  The same is clear for inner
automorphisms defined by sections of \(T\).  Since \(P\) is generated by
\(T\) and by the \(U_\alpha\) with \(a(\alpha)\geq0\), it follows that:
\begin{enumerate}[label=\textup{(\roman*)}]
\item every \(U_i\) is invariant in \(P\);
\item the inner automorphisms defined by sections of \(P\) act linearly on
\(U_i/U_{i+1}\simeq W(\mathcal E_i)\).
\end{enumerate}
\SGARefTarget{sga3:xxvi:item:2:1:1}

\paragraph{\SGARefLink{sga3:xxvi:item:2:1:1}{2.1.2}.}
Let \((T',M',R',\Delta',\ldots)\) now be another pinning of \(G\) adapted
to \(P\).  By \SGARefLink{sga3:xxvi:prop:1:15}{1.15} there is an inner automorphism of \(G\), coming from
\(P\), that carries the first pinning to the second.  After extending
\(S\), it may be written \(\Int(p)\) for some \(p\in P(S)\).
Repeating the construction with the second pinning produces groups
\(U'_i\) and vector-group isomorphisms
\[
 U'_i/U'_{i+1}\simeq W(\mathcal E'_i)
\]
obtained from \(U_i\) and
\(U_i/U_{i+1}\simeq W(\mathcal E_i)\) by transport of structure under
\(\Int(p)\).  The preceding invariance and linearity statements therefore
give \(U'_i=U_i\), and the two vector-space structures on
\[
 U_i/U_{i+1}=U'_i/U'_{i+1}
\]
coincide.

Thus the \(U_i\) and the vector structures on their successive quotients
are independent of the chosen pinning and, in particular, invariant under
every automorphism of \(P\).  This proves the proposition when \((G,P)\)
admits a pinning.  Assertion (iii) is immediate: by Exp.\ XXI, \SGARefLink{sga3:XXI:3.1.2}{3.1.2}, the
set \(\{a(\gamma)\mid\gamma\in R\}\) is an interval in \(\mathbb Z\), so
\(\dim U_{i,s}=\dim U_{i+1,s}\) cannot occur unless \(U_{i,s}=e\).
\SGARefTarget{sga3:xxvi:item:2:1:2}

\paragraph{\SGARefLink{sga3:xxvi:item:2:1:2}{2.1.3}.}
In general there is an fpqc covering family \(\{S_j\to S\}\) such that
every pair \((G_{S_j},P_{S_j})\) admits a pinning (\SGARefLink{sga3:xxvi:lemma:1:14}{1.14}).  What precedes
gives descent data on the \(\rad^u(P_{S_j})_i\), compatible with the
vector structures on the quotients.  The result follows by descent for
closed subschemes and for locally free modules.
\SGARefTarget{sga3:xxvi:item:2:1:3}

\begin{sgastatement}{Corollary 2.2}
Let \(S\) be affine, \(G\) an \(S\)-reductive group, and \(P\) a parabolic
subgroup of \(G\).  Then
\[
 H^1(S,\rad^u(P))=0;
\]
that is, every principal homogeneous bundle under \(\rad^u(P)\) is
trivial.
\end{sgastatement}
\SGARefTarget{sga3:xxvi:cor:2:2}

Indeed, \(S\) is the disjoint union of subschemes on each of which
\(\rad^u(P)\) has constant relative dimension.  By 2.1(iii) one may thus
assume that \(U_n=e\) for some \(n\).  Since
\[
 H^1(S,U_i/U_{i+1})
 =H^1(S,W(\mathcal E_i))=0
\]
by TDTE I, B, \SGARefLink{sga3:xxvi:def:1:1}{1.1} (or SGA~1, \SGARefLink{sga3:xi:corollary:5:1}{Exposé XI, 5.1}), the conclusion follows.

\begin{sgastatement}{Corollary 2.3}
Under the preceding hypotheses, \(P\) has a Levi subgroup \(L\).  If
\(L\) is any Levi subgroup of \(P\), the canonical map
\[
 H^1(S,L)\longrightarrow H^1(S,P)
\]
is bijective \textup{(see the introduction to Exp.\ XXIV for the
definition of \(H^1(S,-)\))}.
\end{sgastatement}
\SGARefTarget{sga3:xxvi:cor:2:3}

The first assertion follows from \SGARefLink{sga3:xxvi:cor:2:2}{2.2} and \SGARefLink{sga3:xxvi:cor:1:9}{1.9}.  The map
\(H^1(S,L)\to H^1(S,P)\) is surjective because
\(P=L\cdot\rad^u(P)\).  For injectivity it is enough to show that
\[
 H^1\bigl(S,\rad^u(P)_Q\bigr)=0
\]
for every principal homogeneous \(L\)-bundle \(Q\), where the subscript
denotes twisting by \(Q\).  One may either repeat the proof of \SGARefLink{sga3:xxvi:cor:2:2}{2.2}, using
the \(L\)-invariance of the vector structures on \(U_i/U_{i+1}\), or
observe that \(\rad^u(P)_Q\) is the unipotent radical of the parabolic
subgroup \(P_Q\) of \(G_Q\) and apply \SGARefLink{sga3:xxvi:cor:2:2}{2.2} to \(P_Q\).

\begin{sgastatement}{Corollary 2.4}
Let \(S\) be a semi-local scheme, \(G\) an \(S\)-reductive group, and
\(P\) a parabolic subgroup of \(G\).  There is a maximal torus \(T\) of
\(G\) contained in \(P\).
\end{sgastatement}
\SGARefTarget{sga3:xxvi:cor:2:4}

Indeed, by \SGARefLink{sga3:xxvi:cor:2:3}{2.3} the group \(P\) has a Levi subgroup \(L\), and it suffices
to know that \(L\) has a maximal torus, which follows from Exp.\ XIV,
\SGARefLink{sga3:xiv:3.20}{3.20}.

\begin{sgastatement}{Corollary 2.5}
Let \(S\) be affine, \(G\) an \(S\)-reductive group, and \(P\) a parabolic
subgroup of \(G\).  There is a locally free \(\mathcal O_S\)-module
\(\mathcal E\) such that \(\rad^u(P)\), as an \(S\)-scheme, is isomorphic
to \(W(\mathcal E)\).
\end{sgastatement}
\SGARefTarget{sga3:xxvi:cor:2:5}

We prove by induction on \(i\) that there is an isomorphism of \(S\)-schemes
\[
 \rad^u(P)/U_i
 \simeq
 W(\mathcal E_0\oplus\mathcal E_1\oplus\cdots\oplus\mathcal E_{i-1}).
\]
This is clear for \(i=0\).  If \(i>0\), then \(\rad^u(P)/U_i\) is a
principal homogeneous bundle over \(\rad^u(P)/U_{i-1}\) under
\[
 (U_{i-1}/U_i)_{\rad^u(P)/U_{i-1}}
 \simeq
 W\!\left(
 \mathcal E_{i-1}\otimes
 \mathcal O_{\rad^u(P)/U_{i-1}}
 \right).
\]
The base is affine, for example by induction, so the bundle is trivial
(TDTE~I or SGA~1, Exposé XI, \emph{loc.\ cit.}).  Hence
\[
 \rad^u(P)/U_i
 \simeq
 \bigl(\rad^u(P)/U_{i-1}\bigr)\times_S(U_{i-1}/U_i),
\]
which completes the induction.

\begin{sgastatement}{Corollary 2.6}
Let \(S\) be semi-local, let \(\{s_i\}\) be its closed points, let \(G\)
be an \(S\)-reductive group, and let \(P\) be a parabolic subgroup of
\(G\).  The canonical map
\[
 \rad^u(P)(S)\longrightarrow
 \prod_i\rad^u(P)(\operatorname{Spec}\kappa(s_i))
\]
is surjective.
\end{sgastatement}
\SGARefTarget{sga3:xxvi:cor:2:6}

Indeed, write \(S=\operatorname{Spec}A\) and
\(\kappa(s_i)=A/\mathfrak p_i\).  If \(\mathcal E\) corresponds to the
projective, hence flat, \(A\)-module \(E\), it remains to prove that
\[
 E\longrightarrow\prod_i E\otimes_A A/\mathfrak p_i
\]
is surjective.  It suffices to do so for \(E=A\), when it is well known
(Bourbaki, \emph{Algèbre commutative}, Chapter~II, §1, no.~2,
Proposition~5).

\begin{sgastatement}{Corollary 2.7}
Let \(k\) be an infinite field, \(G\) a \(k\)-reductive group, and \(P\)
a parabolic subgroup of \(G\).  Then \(\rad^u(P)(k)\) is dense in
\(\rad^u(P)\).
\end{sgastatement}
\SGARefTarget{sga3:xxvi:cor:2:7}

\begin{sgastatement}{Corollary 2.8}
Let \(S\) be semi-local, let \(\{s_i\}\) be its closed points, let \(G\)
be an \(S\)-reductive group, let \(P\) be a parabolic subgroup of \(G\),
and, for every \(i\), let \(L_i\) be a Levi subgroup of \(P_{s_i}\).
There is a Levi subgroup \(L\) of \(P\) whose fiber at every \(s_i\) is
\(L_i\).
\end{sgastatement}
\SGARefTarget{sga3:xxvi:cor:2:8}

Indeed, let \(L_0\) be a Levi subgroup of \(P\) (\SGARefLink{sga3:xxvi:cor:2:3}{2.3}).  For every \(i\),
choose
\[
 u_i\in\rad^u(P)(\operatorname{Spec}\kappa(s_i))
\quad\text{with}\quad
 \Int(u_i)L_{0,s_i}=L_i
\]
by \SGARefLink{sga3:xxvi:cor:1:8}{1.8}.  If \(u\in\rad^u(P)(S)\) induces the \(u_i\), as \SGARefLink{sga3:xxvi:cor:2:6}{2.6} permits,
then \(L=\Int(u)L_0\) has the required fibers.

\begin{sgastatement}{\SGAInlineRefTarget{sga3:xxvi:corollary:2-9}Corollary 2.9}
In the setting of \SGARefLink{sga3:xxvi:proposition:2-1}{2.1}, let \(H\) also be a sub-group-scheme of \(G\),
smooth and of finite presentation over \(S\), with connected fibers, such
that \(P\cap H\) contains, locally for the fpqc topology, a maximal torus
of \(G\).  For every \(i\geq0\) there is a locally direct-factor
submodule \(\mathcal F_i\subset\mathcal E_i\) for which the isomorphism
\[
 U_i/U_{i+1}\xrightarrow{\sim}W(\mathcal E_i)
\]
induces an isomorphism
\[
 (U_i\cap H)/(U_{i+1}\cap H)\xrightarrow{\sim}W(\mathcal F_i).
\]
\end{sgastatement}
\SGARefTarget{sga3:xxvi:cor:2:9}

Indeed, \(H\) is a subgroup of \(G\) of type (R) (Exp.\ XXII, \SGARefLink{sga3:xxii:def:5.2.1}{5.2.1}).
The assertion is fpqc-local, so we may suppose that \(G\) is split
relative to a maximal torus of \(P\cap H\), and even reduce to \SGARefLink{sga3:xxvi:prop:2:1}{2.1.1} with
\(H\) defined by a subset \(R'\subset R\).  In the notation used there,
Exp.\ XXII, \SGARefLink{sga3:xxii:cor:5.6.7:item:ii}{5.6.7(ii)} gives
\[
 U_i\cap H=\prod_{\alpha\in R_i\cap R'}U_\alpha.
\]
Thus \((U_i\cap H)/(U_{i+1}\cap H)\) identifies with
\[
 \prod_{\substack{\alpha\in R'\\a(\alpha)=i+1}}U_\alpha,
\]
which proves the assertion.

\begin{sgastatement}{Corollary 2.10}
In the setting of \SGARefLink{sga3:xxvi:corollary:2-9}{2.9}, the conclusions of \SGARefLink{sga3:xxvi:cor:2:2}{2.2}, \SGARefLink{sga3:xxvi:cor:2:5}{2.5}, \SGARefLink{sga3:xxvi:cor:2:6}{2.6}, and \SGARefLink{sga3:xxvi:cor:2:7}{2.7} remain
valid with \(\rad^u(P)\) replaced by \(\rad^u(P)\cap H\).
\end{sgastatement}
\SGARefTarget{sga3:xxvi:cor:2:10}

\begin{sgastatement}{Corollary 2.11}
Let\footnote{\emph{Editor's note.} \SGARefLink{sga3:xxvi:cor:2:11}{Corollary 2.11} has been added; it is
used in \SGARefLink{sga3:xxvi:proposition:4-5-1}{4.5.1} and \SGARefLink{sga3:xxvi:cor:6:17}{6.17}.} \(G\) be an \(S\)-reductive group, \(P\) a parabolic subgroup, and
\(H\) a subgroup of \(P\) of type \textup{(RC)} such that
\[
 \rad^u(H)=\rad^u(P)\cap H.
\]
Then statements \SGARefLink{sga3:xxvi:cor:2:2}{2.2} through \SGARefLink{sga3:xxvi:cor:2:8}{2.8} remain valid with \(P\) replaced by
\(H\).
\end{sgastatement}
\SGARefTarget{sga3:xxvi:cor:2:11}
